\section{Main Theorem}

\begin{theorem}[Anchored coefficient-normalized Pfaffian sweep]
\label{thm:main}
Under Assumptions~\ref{assump:parameter-regime},
\ref{assump:balcan-common-chain}, and
\ref{assump:anchored-unit-range}, the following conclusions hold.  The
probabilistic conclusions additionally use
Assumption~\ref{assump:cube-density-laws}, and the general affine conclusion
also uses Assumption~\ref{assump:affine-chart-data}.

\emph{Primitive presentation and projective speed.}
Define
\[
D_*:=\Delta B_Q(1+qB_P).
\]
For every $x\in[-1,1]$ and $1\leq i\leq N$,
\[
|G_i'(x)|\leq D_*,
\qquad
\|G'(x)\|_2\leq\sqrt N\,D_*.
\]
Moreover $G_1=F_1=1$, so $\|G(x)\|_2,\|F(\theta)\|_2\geq1$.  With
$\gamma_G=G/\|G\|_2$,
\[
\gamma_G'(x)
=\frac{(I_N-\gamma_G(x)\gamma_G(x)^{\mathsf T})G'(x)}{\|G(x)\|_2},
\]
\[
\gamma_F(\theta)=\gamma_G\!\left(\frac{\theta-c}{h}\right),
\qquad
\gamma_F'(\theta)=\frac1h\gamma_G'\!\left(\frac{\theta-c}{h}\right),
\]
and hence
\[
\Gamma_{\mathrm{proj}}(F)
\leq\frac{\sqrt N\,\Delta B_Q(1+qB_P)}{h}.
\tag{R1}
\]
These pointwise statements include one-sided endpoint derivatives.  When
$q=0$, $M=B_P=0$ and $D_*=\Delta B_Q$; dependence on $M$ is degree zero
when $B_P$ is fixed.  When $N=1$, or when the Euclidean projector annihilates
$G'$, the normalized curve is stationary and its speed is zero.  The case
$\Delta=0$ likewise gives the identity $0\leq0$.

\emph{Central swept-hyperplane bound.}
Fix the deterministic presentation before selecting a law or interval.  For
every, possibly correlated, $\mu\in\mathcal D_{N,R,\kappa}$ and then every
positive-length interval $I\subseteq\Theta$,
\[
\Pr_{\alpha\sim\mu}
 [\exists\theta\in I:\langle\alpha,F(\theta)\rangle=0]
\leq A\sqrt{\frac N2}\,\Gamma_{\mathrm{proj}}(F)|I|
\leq\frac{AN\Delta B_Q(1+qB_P)}{\sqrt2\,h}|I|,
\tag{R2a}
\]
and, with the interval supremum inside the law supremum,
\[
C^{\mathrm{Pf}}_{\mathcal D}(F;\Theta)
\leq A\sqrt{\frac N2}\,\Gamma_{\mathrm{proj}}(F)
\leq\frac{AN\Delta B_Q(1+qB_P)}{\sqrt2\,h}.
\tag{R2b}
\]
Before division by $|I|$, the same chain holds for empty and singleton
intervals with zero right side and for every endpoint convention.  Tangent
and multiple roots, stationary normalized pieces, and coefficients for which
the combination vanishes identically are included; no simple-root or finite-
fiber condition is imposed.

\emph{General affine chart bound.}
For every interval $I\subseteq\Theta$, every measurable pivot partition in
Assumption~\ref{assump:affine-chart-data}, and every, possibly correlated,
$\mu\in\mathcal D_{N,R,\kappa}$,
\[
\begin{aligned}
&\Pr_{\alpha\sim\mu}
 [\exists\theta\in I:F_0(\theta)+\langle\alpha,F(\theta)\rangle=0]\\
&\quad\leq
\kappa\sum_{j=1}^N\int_{E_j}\int_{[-R,R]^{N-1}}
 |\partial_\theta T_j(\theta,\beta)|\,d\beta\,d\theta.
\end{aligned}
\tag{R3}
\]
This is an inequality in $[0,+\infty]$.  It permits a divergent chart
integral, arbitrarily small nonzero pivots, cube faces, empty or singleton
intervals, included endpoints, tangent and multiple roots, infinite fibers,
$N=1$, and affine combinations that vanish identically.  It assumes neither
transversality, a uniform pivot margin, a conditional density, nor coordinate
independence.

\emph{Exact affine-monic recovery.}
For every integer $d\geq1$ and bounded interval $J\subset\mathbb R$, choose a
nondegenerate $\Theta=[c-h,c+h]\supseteq J$ and set
\[
F_0(\theta)=\theta^d,
\qquad
F_{k+1}(\theta)=\theta^k\quad(0\leq k\leq d-1),
\]
\[
Q_0(x)=(c+hx)^d,
\qquad
Q_{k+1}(x)=(c+hx)^k\quad(0\leq k\leq d-1).
\]
Then
\[
p_\alpha(\theta)
:=F_0(\theta)+\langle\alpha,F(\theta)\rangle
=\theta^d+\sum_{k=0}^{d-1}\alpha_k\theta^k,
\]
with the exact data
\[
q=0,
\quad M=0,
\quad B_P=0,
\quad N=d,
\quad A=(2R)^d\kappa,
\quad \Delta_{\mathrm{aug}}=d.
\]
The monic coefficient $1$ is deterministic and remains outside the random
lower-coefficient vector $\alpha=(\alpha_0,\ldots,\alpha_{d-1})$.

For $d\geq2$, use
\[
E_1=J\cap\{|\theta|\leq1\},
\qquad
E_d=J\cap\{|\theta|>1\},
\qquad
E_j=\varnothing\quad(j\notin\{1,d\}).
\]
The transition points $0,1,-1$ belong to the constant-pivot cell.  The two
charts and their literal velocities are
\[
T_1(\theta,\beta)
=-\theta^d-\sum_{k=1}^{d-1}\beta_k\theta^k,
\qquad
|\partial_\theta T_1|
\leq d+R\sum_{k=1}^{d-1}k
=d+\frac{Rd(d-1)}2
\quad(|\theta|\leq1),
\]
\[
T_d(\theta,\beta)
=-\theta-\sum_{k=0}^{d-2}\beta_k\theta^{k-d+1},
\]
\[
|\partial_\theta T_d|
\leq1+R\sum_{k=0}^{d-2}(d-1-k)
=1+\frac{Rd(d-1)}2
\leq d+\frac{Rd(d-1)}2
\quad(|\theta|>1).
\]
For $d=1$, use $E_1=J$, the zero-dimensional beta cube,
$T_1(\theta)=-\theta$, and $|T_1'(\theta)|=1$.  Consequently, for every
possibly correlated $\mu\in\mathcal D_{d,R,\kappa}$ and then every bounded
$J$,
\[
\Pr_{\alpha\sim\mu}[\exists\theta\in J:p_\alpha(\theta)=0]
\leq\kappa(2R)^{d-1}
 \left(d+\frac{Rd(d-1)}2\right)|J|.
\tag{R4}
\]
At $d=1$ this is exactly $\kappa|J|$.  There is no chart-count factor,
interval enlargement, law augmentation, conservative remainder, or singular
leading coordinate.

\emph{Counter-example 1 scale.}
Let $0<\delta\leq1$, $\Theta=[-1,1]$, $c=0$, $h=1$, and
\[
G(x)=\left(1,\frac{x}{\delta}\right),
\qquad
F(\theta)=\left(1,\frac{\theta}{\delta}\right).
\]
For the uniform law on $[-1,1]^2$,
\[
(h,q,M,\Delta,N,R,\kappa,A,B_P,B_Q)
=\left(1,0,0,1,2,1,\frac14,1,0,\frac1\delta\right),
\]
\[
\gamma_F(\theta)
=\frac{(\delta,\theta)}{\sqrt{\delta^2+\theta^2}},
\qquad
\|\gamma_F'(\theta)\|_2
=\frac{\delta}{\delta^2+\theta^2},
\qquad
\Gamma_{\mathrm{proj}}(F)=\frac1\delta.
\]
For every $0<\epsilon\leq\delta$,
\[
\Pr[\exists\theta\in[0,\epsilon]:
 \alpha_1+\alpha_2\theta/\delta=0]
=\frac{\epsilon}{4\delta}.
\]
Therefore, for $\mathcal D=\mathcal D_{2,1,1/4}$,
\[
\frac1{4\delta}
\leq C^{\mathrm{Pf}}_{\mathcal D}(F;[-1,1])
\leq\underbrace{A\sqrt{\frac N2}\,\Gamma_{\mathrm{proj}}(F)}_{1/\delta}
\leq\underbrace{\frac{AN\Delta B_Q(1+qB_P)}{\sqrt2\,h}}_{\sqrt2/\delta}.
\tag{R5}
\]
This is a lower and upper scale certificate, not an equality or optimality
claim for the capacity.  It includes $\delta=1$ and $\epsilon=\delta$
without a limiting argument.

\emph{Rate contract.}
The exposed central variables are
$(q,M,\Delta,N,R,\kappa,A,B_P,B_Q,h^{-1})$; the monic variables are
$(d,R,\kappa,|J|)$; and the scale-audit variables are $(\delta,\epsilon)$.
There are no hidden constants or allowed hidden parameter dependencies.  The
deterministic presentation is fixed before the law and interval ranges.  The
probability mode is ordinary probability under each fixed full joint law; the
horizon mode is the stated per-interval assertion followed, only in
\emph{(R2b)}, by the interval supremum and then the law supremum.  Central
norms and operator norms are Euclidean, affine velocities are scalar, and all
coefficient and interval measures use the conventions in
Section~\ref{sec:prelim}.  The
admissibility conditions are exactly the numbered assumptions and the local
restrictions $d\geq1$, bounded $J$, and
$0<\epsilon\leq\delta\leq1$.  No confidence parameter, asymptotic horizon,
independence premise, or further normalization claim is present, and every
displayed algebraic bridge and numerical factor is literal.
\end{theorem}
